% Vietnamese translation for OI Public Library
% Source-File: IOI中国国家候选队论文/2008/Day1/6.肖汉骏《例谈信息学竞赛分析中的“深”与“广”》/讲稿.doc
\documentclass[11pt]{article}
\usepackage{oipl-vn}

\newcommand{\Oh}{\mathcal{O}}

\title{Đường rẽ còn dài, trên dưới kiếm tìm\\{\large Bàn về ``chiều sâu'' và ``độ rộng'' trong phân tích bài toán tin học}}
\author{Xiao Hanjun\\Trường Trung học số 1 Phúc Châu}
\date{}

\begin{document}
\maketitle

\origtitle{Depth and Breadth in Informatics Competition Analysis}
\sourcepath{IOI 2008 / Day 1 / Xiao Hanjun / talk-script.doc}

\section{Mở đầu}

Xin chào mọi người, tôi là Xiao Hanjun từ Trường Trung học số 1 Phúc Châu.
Chủ đề tôi trình bày hôm nay là: \emph{Đường rẽ còn dài, trên dưới kiếm tìm:
bàn về ``chiều sâu'' và ``độ rộng'' trong phân tích bài toán tin học}.

Người ta thường dùng ``sâu'' và ``rộng'' để mô tả đại dương. Nhưng ở một góc
nhìn khác, trí óc con người còn sâu và rộng hơn bất kỳ đại dương nào. Trong
một sân khấu nơi năng lực tư duy được bộc lộ đầy đủ như thi tin học,
``chiều sâu'' và ``độ rộng'' có ý nghĩa gì? Trong thực hành giải bài, làm thế
nào để đạt được cả hai điều ấy? Bài nói này cùng thảo luận các câu hỏi đó.

\section{Chiều sâu và độ rộng}

Trong phân tích bài toán, ``chiều sâu'' có mấy tầng nghĩa.

Thứ nhất là \term{tính phân tầng}. Quan sát theo tầng nghĩa là nắm được mạch
nguồn của bài toán, bắt đầu từ trường hợp ít chiều hơn, điều kiện đơn giản
hơn, rồi đi từ nông tới sâu để phân tích sự thay đổi khi bài toán lên chiều
cao hơn hoặc điều kiện phức tạp hơn.

Thứ hai là \term{tính liên tục}. Quá trình phân tích thường dài và gập ghềnh.
Nếu nghĩ đến đâu phân tích đến đó, tư duy rất dễ rối. Ngược lại, nếu có thể
đi sâu từng bước theo một hướng xác định và tận dụng kết quả của mỗi bước,
ta dễ đào ra những thông tin hữu ích hơn.

Thứ ba là chú ý đến \term{quan hệ giữa các yếu tố}. Bài toán tin học không
phải chỉ là tổng đơn giản của các yếu tố. Quan hệ giữa chúng cũng tồn tại một
cách khách quan, thậm chí bản chất và sâu sắc hơn. Những quan hệ như tính đơn
điệu hoặc chu kỳ thường có thể trở thành điểm đột phá.

Còn ``độ rộng'' trong phân tích bài toán có nội hàm phong phú hơn.

Thứ nhất là \term{tầm nhìn rộng}. Trong thời gian gấp gáp của phòng thi, có
nhiều bài toán khó, thậm chí không thể, được giải hoàn hảo ngay. Khi đó ta
cần phát huy trí tưởng tượng và dùng nhiều chiến lược khác nhau.

Thứ hai là \term{mạch suy nghĩ rộng}. Cùng một bài toán, nhìn từ các góc khác
nhau thường cho các kết quả khác nhau. Nếu tổng hợp được các kết quả ấy, ta
sẽ hiểu bài toán toàn diện hơn; những thông tin ẩn cũng khó thoát khỏi các
hướng thăm dò khác nhau.

Thứ ba là \term{phương tiện phân tích phong phú}. Phương tiện phân tích là
công cụ sắc bén để nắm bản chất bài toán. Nếu nắm được nhiều phương tiện, ta
giống như người luyện đủ nhiều loại vũ khí: trong lúc nói cười, bài toán đã
tan biến.

\section{Ví dụ: thứ tự đọc bài viết}

Xét một bài toán: có \(N\) bài viết. Bài thứ \(i\) có thời gian đọc \(T_i\)
và giá trị \(V_i\). Chi phí đọc một bài được định nghĩa là thời điểm đọc xong
bài đó nhân với giá trị của nó. Cần tìm thứ tự đọc sao cho tổng chi phí nhỏ
nhất.

Ngoài ra, một số bài có cùng tác giả; các bài của cùng một tác giả bắt buộc
phải được đọc theo đúng thứ tự thời gian mà tác giả viết ra.

Bài toán yêu cầu tìm một thứ tự tối ưu nhưng kèm thêm ràng buộc theo tác giả.
Xét trực tiếp bài toán đầy đủ không dễ tìm hướng tốt, nên ta trước hết làm
yếu điều kiện và giải trường hợp đặc biệt: bỏ qua yếu tố tác giả.

Trong trường hợp đặc biệt, ta cần tìm một hoán vị tối ưu. Hoán vị tối ưu chắc
chắn không phải được sắp tùy tiện; điểm đặc biệt của nó là mọi thay đổi cục
bộ đều không làm kết quả tốt hơn. Thay đổi đơn giản nhất trong một dãy là đổi
chỗ hai phần tử kề nhau.

Xét hai bài kề nhau có thời gian \(T_1,T_2\) và giá trị \(V_1,V_2\). Nếu đổi
chỗ chúng, bài thứ nhất bị hoàn thành muộn thêm \(T_2\), còn bài thứ hai
được hoàn thành sớm hơn \(T_1\). Biến thiên của tổng chi phí là
\[
  T_2V_1-T_1V_2.
\]
Với thứ tự tối ưu, biến thiên này phải dương, tức
\[
  T_2V_1>T_1V_2,
\]
hay
\[
  \frac{V_1}{T_1}>\frac{V_2}{T_2}.
\]
Nếu mọi cặp kề nhau đều thỏa quan hệ này, toàn bộ dãy phải được sắp theo
\(V_i/T_i\) giảm dần. Trường hợp đặc biệt được giải trong
\(\Oh(N\log N)\).

\section{Trở lại ràng buộc tác giả}

Ta quay lại bài toán tổng quát. Ảnh hưởng của yếu tố tác giả nên xử lý thế
nào? Vẫn áp dụng tư tưởng từ đặc biệt đến tổng quát. Trong mọi trường hợp
tổng quát, trường hợp có hai tác giả là đặc biệt nhất; trong các trường hợp
hai tác giả, trường hợp một tác giả chỉ viết một bài lại đặc biệt nhất. Hãy
xét xem bài độc lập ấy nên được đọc khi nào.

Theo phân tích trên, trong thứ tự tối ưu, các phần tử kề nhau có thể đổi chỗ
phải thỏa
\[
  \frac{V_1}{T_1}>\frac{V_2}{T_2}.
\]
Nói cách khác, tỉ số của bài độc lập phải nhỏ hơn bài phía trước và lớn hơn
bài phía sau. Nhưng chỉ suy diễn đại số thì có thể có nhiều vị trí thỏa điều
kiện, không có tính chất đủ tốt để khai thác.

Dạng tỉ số của cặp \((V,T)\) gợi đến hệ số góc của đường thẳng. Ta thử dùng
một công cụ phân tích rộng hơn: hình học. Xem mỗi bài viết là một vectơ
\[
  (V_i,T_i),
\]
thì \(V_i/T_i\) chính là hệ số góc.

Kết luận đại số ở trên trở thành: hệ số góc của bài độc lập phải nằm giữa hệ
số góc của hai phía trước sau. Vị trí chèn bài độc lập có thể được biểu diễn
trực quan như một \term{điểm lồi}.

Nhưng nếu chỉ chú ý đến một mặt của bài toán thì trái với yêu cầu về độ rộng.
Hãy nhìn cả các vị trí không thể chèn. Trong hình học, chúng tương ứng với
\term{điểm lõm}. Có thể dùng phản chứng để chứng minh rằng điểm lõm không thể
được chèn thêm bài mới. Vì thế ta có thể gộp hai vectơ tạo thành điểm lõm để
loại bỏ nó.

Hình nào không chứa điểm lõm? Đó chính là bao lồi. Đến đây, ta chỉ còn cách
đáp án một bước.

\section{Thuật toán}

Với mỗi tác giả, xử lý riêng các bài của tác giả đó và dựng bao lồi của các
vectơ bài viết. Những điểm trên bao lồi chia dãy bài của tác giả thành một số
đoạn. Sau đó, lấy tất cả các đoạn bài viết và sắp chúng theo hệ số góc giảm
dần; thứ tự thu được là thứ tự tối ưu.

Nhìn lại quá trình giải bài toán này, ta thấy các bước chính là:
\begin{itemize}
  \item trước hết bỏ ràng buộc, xét trường hợp đặc biệt, tìm được quan hệ
  giữa hai phần tử kề nhau trong thứ tự tối ưu;
  \item trong bài toán tổng quát, tiếp tục chọn trường hợp hai tác giả và một
  bài độc lập để phân tích;
  \item dùng kết hợp số và hình, xét cả điểm lồi lẫn điểm lõm;
  \item cuối cùng liên hệ đến bao lồi và giải trọn bài toán.
\end{itemize}
Mỗi bước trong quá trình ấy đều ít nhiều gắn với ``chiều sâu'' và ``độ rộng''.

\section{Kết luận}

Trong thực hành giải bài hằng ngày, chúng ta thường dựa vào kinh nghiệm và
vô thức dùng nhiều phương pháp phân tích. Hiệu quả khi tốt khi xấu, khó nắm
bắt. Điều này cho thấy kinh nghiệm rời rạc không thể thay thế lý thuyết có
hệ thống.

Nếu trong quá trình phân tích ta tuân theo một số quy tắc, phân tích có lý và
có thứ tự, ta có thể tránh sự bất ổn khi chỉ dựa vào kinh nghiệm. Hơn nữa,
giải được bài toán chưa phải điểm kết thúc. Trong luận văn, tác giả còn trình
bày cách mở rộng bài toán theo hướng ``sâu'' và ``rộng'' để nâng cao năng lực
phân tích.

Mài dao không làm chậm việc chặt củi. Nếu khi giải bài ta biết định mưu rồi
mới hành động, kiên trì thói quen phân tích tốt; sau khi giải xong còn có ý
thức mở rộng bài toán theo chiều sâu và độ rộng, so sánh lời giải ở các tầng
và góc nhìn khác nhau, biết cả kết quả lẫn nguyên nhân của nó, thì theo thời
gian năng lực tư duy và phân tích sẽ tiến bộ rõ rệt. Khi đó ta mới thật sự
đạt điều mà Olympic theo đuổi: phát triển trí tuệ và nâng cao năng lực.

\end{document}
